\section{المصطلحات}
\label{sec:terminology}

تميّز المصطلحات الآتية بين قناة التواصل العامة، والإعداد المشترك لدى الطرفين، والأدلة المستخدمة لتقييم استخراج الرسالة.

\begin{table}[htbp]
\centering\small
\caption{المصطلحات الأساسية المستخدمة في وصف الطريقة.}
\label{tab:terminology}
\begin{RTL}
\begin{tabularx}{\linewidth}{|A{0.25\linewidth}|A{0.25\linewidth}|B|}
\hline
المصطلح العربي & المقابل الإنجليزي & المقصود في الطريقة \\
\hline
الحمولة & \EN{Payload} & الرسالة السرية المراد نقلها. \\
\hline
النص الحامل & \EN{Cover text} & النص الطبيعي المستخدم لحمل المعلومات الخفية. \\
\hline
الإخفاء باستخدام الصندوق الأسود & \EN{Black-box steganography} & استخدام نموذج الكتابة عبر واجهته المتاحة. \\
\hline
\EN{Prompt} & \EN{Prompt} & التعليمات والسياق المقدمان إلى النموذج. \\
\hline
النموذج المساعد الحتمي & \EN{Deterministic helper LLM} & نموذج بإعداد ثابت لبناء الاستعلامات والموضوعات. \\
\hline
نموذج الكتابة & \EN{Writer LLM} & النموذج الذي يصوغ التعليق الحامل. \\
\hline
نسخة من تعليقات المنشور & \EN{Comments snapshot / dump} & التعليقات التي يعتمد عليها الطرفان لبناء القوائم المشتركة. \\
\hline
قائمة النصوص المصدرية & \EN{Source-text list} & نصوص المنشور والتعليقات ووثائق المجال، بترتيب مشترك. \\
\hline
قائمة موضوعات التعليق & \EN{Comment-topic list} & موضوعات مرتبة، مع اقتباس وفئة ومرجع مصدر لكل موضوع. \\
\hline
موضع الرد & \EN{Reply location} & المنشور أو التعليق الذي يرد عليه المُرسِل. \\
\hline
قاموس الضغط & \EN{Compression dictionary} & مواد مشتركة تمثل مقاطع الحمولة بمؤشرات أقصر. \\
\hline
الإطار & \EN{Frame} & تعليق حامل ضمن رسالة موزعة على عدة تعليقات. \\
\hline
\end{tabularx}
\end{RTL}
\end{table}
\FloatBarrier

\paragraph{الحمولة.}
الرسالة السرية المراد إرسالها. وتختلف بايتات الرسالة الأصلية عن تمثيلها المضغوط وعن حقول النقل المضافة إليها.

\paragraph{النص الحامل والنص المتضمن للرسالة.}
النص الحامل نص عادي يُستخدم لنقل المعلومات الخفية. أما النص المتضمن للرسالة فهو النص المنشور الذي يحمل هذه المعلومات. وفي طريقتنا، يتكون الحامل من التعليق المولَّد وموضع الرد معًا.

\paragraph{خفاء الرسالة المضمَّنة.}
صعوبة ملاحظة وجود التواصل الخفي. ويختص الجانب الإدراكي بحكم القارئ على طبيعية النص، والجانب الإحصائي باختلاف توزيعات النصوص، والجانب المعرفي باتساق المعنى وملاءمة السياق.

\paragraph{نموذج اللغة الكبير والرمز.}
يولّد نموذج اللغة الكبير، واختصاره \EN{LLM}، النص بتوقع الرموز المتتابعة. والرمز وحدة من مفردات النموذج، وليس بالضرورة كلمة كاملة.

\paragraph{\EN{Prompt}.}
التعليمات والسياق المقدَّمان إلى نموذج اللغة. ويتضمن \EN{prompt} كتابة التعليق المنشورَ، وسلسلة التعليقات الأب ذات الصلة، وموضوع التعليق المختار.

\paragraph{الإخفاء باستخدام الصندوق الأسود.}
استخدام النموذج عبر واجهته المتاحة دون الاطلاع على احتمالات الرموز داخله. وفي طريقتنا، تحدد البتات موضع الرد وموضوعه قبل أن يكتب النموذج التعليق.

\paragraph{التوليد دون تدريب مخصص.}
توليد النص دون إجراء \EN{fine-tuning} للنموذج المدرَّب مسبقًا على بيانات المهمة. ويمكن مع ذلك توجيه التوليد بأمثلة ضمن \EN{prompt}. ولا تحدد هذه الخاصية وحدها ما إذا كانت الطريقة تتطلب الوصول إلى النموذج بوصفه صندوقًا أسود أو أبيض.

\paragraph{مَوْضِع الرَّدّ وسلسلة التعليقات الأب.}
موضع الرد هو المنشور أو التعليق الموجود الذي يرد عليه المُرسِل. وتتكون سلسلة التعليقات الأب من التعليقات السابقة التي يصل عبرها ذلك التعليق إلى المنشور، وتوفر سياق المحادثة المحلي.

\paragraph{مَوْضُوع التَّعْلِيق وقائمة الموضوعات.}
يحدد موضوع التعليق ما يتحدث عنه الرد المولَّد. وتربط القائمة المرتبة الموضوعات المرشحة بفهارسها وموادها المصدرية. ويحتوي كل موضوع على اقتباس مصدر ووصف وفئة ومرجع للنص الذي ألهمه. ونستخدم اسم قائمة موضوعات التعليق للبنية المسماة \EN{tangent DB} في التنفيذ.

\paragraph{قناة الاختيار.}
القناة التي تحدد فيها بتات الرسالة موضع الرد وموضوع التعليق، ويستخرج المُستقبِل البتات من فهرسيهما. وتحمل هذه الاختيارات البتات دون ربطها باختيارات الرموز المفردة.

\paragraph{نسخة من تعليقات المنشور ونافذة المزامنة.}
تضم نسخة من تعليقات المنشور التعليقات التي يعتمد عليها الطرفان مع علاقات الرد بينها. وتحدد النافذة المتفق عليها $T_{sync}$ زمن قطع التعليقات المقصود بطرح $T_{sync}$ من وقت النشر. ويختلف تطبيق هذا القطع عن إعادة الحصول على المواد المصدرية نفسها.

\paragraph{قائمة النصوص المصدرية.}
قائمة مرتبة تجمع عنوان المنشور ونصه، ونسخة من تعليقات المنشور، ووثائق المجال. يحولها النموذج المساعد إلى قائمة موضوعات التعليق وفق إجراء وإعداد مشتركين.

\paragraph{الإعداد المشترك.}
السياق، وقائمة الموضوعات المرتبة، وقاموس الضغط، ومعاملات التشغيل التي يحتاج الطرفان إلى إعادة بنائها أو الحصول عليها وفق شروط التجربة. وتكشف مقارنة التجزئات الاختلافات، لكنها لا تعيد بناء المواد المفقودة.

\paragraph{قاموس الضغط.}
نص مشترك يستخدمه الضاغط لاستبدال المقاطع المطابقة من الحمولة بمؤشرات. ويختلف هذا القاموس عن قائمة الموضوعات المستخدمة لاختيار موضوع التعليق.

\paragraph{التحقق لدى المُرسِل.}
اختبار المُرسِل أن التعليق المرشح يعطي الموضوع المختار قبل النشر. ولا يثبت اجتياز هذا الاختبار أن مُستقبِلًا مستقلًا سيحدد الموضوع نفسه لاحقًا.

\paragraph{إعادة البناء المدعومة بالتدقيق واستخراج الرسالة لدى المُستقبِل.}
تستخدم إعادة البناء المدعومة بالتدقيق سجلات المُرسِل الخاصة، مثل الفهارس المختارة أو البتات غير المضمَّنة، لإعادة بناء الرسالة. أما استخراج الرسالة لدى المُستقبِل فيعتمد على الحوامل المرئية والمدخلات المشتركة أو المزودة المحددة للتجربة، دون تلك السجلات.

\paragraph{الإطار.}
تعليق حامل واحد ضمن إرسال يتكون من عدة تعليقات. وتتضمن الرسالة المؤطرة أيضًا معلومات تحكم للطول وعدد الإطارات والحشو وفحوص السلامة. وتستهلك هذه الحقول جزءًا من ميزانية البتات المنقولة.

\paragraph{ميزانية الاختيار بعرض ثابت.}
عند توفر $C$ موضعًا للرد و$A$ موضوعًا، تتيح مجموعتان مستقلتان، حجم كل منهما قوة للعدد اثنين، $\lfloor\log_2 C\rfloor+\lfloor\log_2 A\rfloor$ بتات قابلة للتمييز لكل تعليق. وتختلف هذه الميزانية المتاحة عن الحمولة المضمَّنة فعليًا وعن بتات الرسالة المفيدة التي يستخرجها المُستقبِل مستقلًا.

\paragraph{\EN{Perplexity}.}
مقياس لاحتمال النص وفق نموذج لغة مستخدم للتقييم. وتشير القيمة الأقل إلى أن النص أكثر قابلية للتوقع لدى ذلك النموذج، ولا تثبت خفاء الرسالة المضمَّنة أو نجاح استخراجها.
